\section{Предельные теоремы теории вероятностей}

Закон больших чисел (ЗБЧ) математически обосновывает тот факт, что при большом числе случайных экспериментов их средний результат перестает быть случайным и стремится к определенному постоянному значению. Для доказательства ЗБЧ используются неравенства Маркова и Чебышева.

\subsection{Неравенство Маркова}
\begin{theorem}[Неравенство Маркова]
Пусть $X \ge 0$ — неотрицательная дискретная случайная величина с конечным математическим ожиданием $E(X)$. Тогда для любого положительного числа $\varepsilon > 0$:
$$P(X \ge \varepsilon) \le \frac{E(X)}{\varepsilon}.$$
\end{theorem}

\begin{proof}
Запишем определение математического ожидания для дискретной случайной величины. Так как $X \ge 0$, все ее значения $x_i \ge 0$.
Разделим сумму на две части: значения, которые меньше $\varepsilon$, и значения, которые больше или равны $\varepsilon$.
$$E(X) = \sum_{x_i < \varepsilon} x_i P(X = x_i) + \sum_{x_i \ge \varepsilon} x_i P(X = x_i).$$
Первая сумма неотрицательна. Если мы ее отбросим, мы только уменьшим общее выражение:
$$E(X) \ge \sum_{x_i \ge \varepsilon} x_i P(X = x_i).$$
В оставшейся сумме каждый элемент $x_i \ge \varepsilon$. Если мы заменим каждый $x_i$ на меньшее число $\varepsilon$, сумма снова уменьшится:
$$E(X) \ge \varepsilon \sum_{x_i \ge \varepsilon} P(X = x_i) = \varepsilon \cdot P(X \ge \varepsilon).$$
Разделив обе части на $\varepsilon > 0$, получаем искомое неравенство: $P(X \ge \varepsilon) \le \frac{E(X)}{\varepsilon}$.
\end{proof}

\subsection{Неравенство Чебышева}
\begin{theorem}[Неравенство Чебышева]
Для любой случайной величины $X$ с конечными математическим ожиданием $\mu = E(X)$ и дисперсией $D(X)$, для любого $\varepsilon > 0$ справедливо:
$$P(\vert{}X - \mu\vert{} \ge \varepsilon) \le \frac{D(X)}{\varepsilon^2}.$$
\end{theorem}

\begin{proof}
Рассмотрим новую случайную величину $Y = (X - \mu)^2$. Эта величина неотрицательна, а ее математическое ожидание по определению равно дисперсии: $E(Y) = D(X)$.
Событие $\vert{}X - \mu\vert{} \ge \varepsilon$ эквивалентно событию $(X - \mu)^2 \ge \varepsilon^2$, то есть $Y \ge \varepsilon^2$.
Применим доказанное ранее неравенство Маркова для величины $Y$ и константы $\varepsilon^2$:
$$P(Y \ge \varepsilon^2) \le \frac{E(Y)}{\varepsilon^2} \implies P(\vert{}X - \mu\vert{} \ge \varepsilon) \le \frac{D(X)}{\varepsilon^2}.$$
\end{proof}

\subsection{Схема Бернулли и Биномиальное распределение}
Число успехов $S_n$ в $n$ независимых испытаниях Бернулли подчиняется биномиальному закону с вероятностью:
$$P(S_n = k) = C_n^k p^k q^{n-k}.$$
Ее основные числовые характеристики (математическое ожидание и дисперсия) выражаются формулами:
$$E(S_n) = n p, \quad D(S_n) = n p q.$$

\subsection{Законы больших чисел}

\begin{theorem}[Закон больших чисел в форме Бернулли]
Относительная частота $\frac{S_n}{n}$ сходится по вероятности к вероятности успеха $p$:
$$P\left(\left\vert\frac{S_n}{n} - p\right\vert \ge \varepsilon\right) \le \frac{p q}{n \varepsilon^2} \xrightarrow[n \to \infty]{} 0.$$
\end{theorem}

\begin{theorem}[Теорема Чебышева]
Пусть $X_1, X_2, \dots, X_n$ — последовательность попарно независимых случайных величин с математическими ожиданиями $E(X_i) = \mu_i$ и ограниченными дисперсиями $D(X_i) \le C$. Тогда среднее арифметическое этих величин сходится по вероятности к среднему их математических ожиданий:
$$\lim_{n \to \infty} P\left(\left\vert{}\frac{1}{n}\sum_{i=1}^n X_i - \frac{1}{n}\sum_{i=1}^n \mu_i\right\vert{} \ge \varepsilon\right) = 0.$$
\end{theorem}

\begin{proof}
Рассмотрим случайную величину $Z_n = \frac{1}{n}\sum_{i=1}^n X_i$.
По свойствам линейности (см. Теорему \ref{thm:ev_properties} на стр. \pageref{thm:ev_properties}), ее математическое ожидание равно $E(Z_n) = \frac{1}{n}\sum_{i=1}^n \mu_i$.
Поскольку $X_i$ независимы, дисперсия суммы равна сумме дисперсий (согласно Теореме \ref{thm:var_properties}):
$$D(Z_n) = D\left(\frac{1}{n}\sum_{i=1}^n X_i\right) = \frac{1}{n^2} \sum_{i=1}^n D(X_i).$$
Так как каждая дисперсия ограничена константой $C$, то:
$$D(Z_n) \le \frac{1}{n^2} \sum_{i=1}^n C = \frac{1}{n^2} (n \cdot C) = \frac{C}{n}.$$
Применим неравенство Чебышева к величине $Z_n$:
$$P(\vert{}Z_n - E(Z_n)\vert{} \ge \varepsilon) \le \frac{D(Z_n)}{\varepsilon^2} \le \frac{C}{n \varepsilon^2}.$$
При стремлении $n \to \infty$ правая часть стремится к нулю. Следовательно, вероятность отклонения от математического ожидания тоже стремится к нулю.
\end{proof}